Una empresa minera extrae \textbf{100 toneladas} de mineral rojo y \textbf{80 toneladas} de mineral negro cada semana. Estos pueden ser tratados de diferentes maneras para producir tres diferentes aleaciones, blanda, dura o fuerte. Para producir 1 tonelada de aleación blanda se requiere \textbf{5 toneladas} de mineral rojo y \textbf{1 tonelada} de negro. Para la aleación dura los requisitos son \textbf{3 toneladas} de rojo y \textbf{5 toneladas} de negro, mientras que para la aleación fuerte son \textbf{5 toneladas} de rojo y \textbf{3 toneladas} de negro. La ganancia por tonelada de la venta de las aleaciones (después de considerar los costes de producción, pero no los de extracción, que se consideran fijos) son \textbf{250}, \textbf{300} y \textbf{400 euros} por tonelada de aleación blanda, dura y fuerte, respectivamente.

Además, el mineral rojo se puede vender a \textbf{100} euros la tonelada y el mineral negro se debe emplear al completo.

Existe un compromiso comercial de servir al menos \textbf{20} toneladas conjuntamente entre las aleaciones dura y fuerte.

El siguiente modelo de programación lineal permite obtener el mejor plan de explotación de la mina.

\begin{equation}
\begin{split}
\mbox{max. } z = 250x_1 + 300x_2 + 400x_3 + 100x_4\\
s.a.:\\
5x_1 + 3x_2 + 5x_3 + x_4 \leq 100 \hspace{1cm} \mbox{(mineral rojo)} \\
 x_1 + 5x_2 + 3x_3   = 80\hspace{1cm} \mbox{(mineral negro)} \\
        x_2 +  x_3 +  \geq 20 \hspace{1cm} \mbox{(compromiso comercial)} \\
x_1,\,\,x_2,\,\,x_3,\,\,x_4\geq 0\\
\end{split}
\end{equation}

Donde $x_1$, $x_2$ y $x_3$ representan las producciones de aleación blanda, dura y fuerte (en toneladas), respectivamente y $x_4$ son las toneladas de mineral rojo de venta en el mercado.

La siguiente es una tabla correspondiente a una solución intermedia correspondiente a la aplicación del método de las dos fases.

\begin{center}
\begin{tabular}{c|c|cccccccc|}
 & $z$ & $x_1$ & $x_2$ & $x_3$ & $x_4$ & $h_1$ & $h_3$ & $a_2$ & $a_3$\\ 
Fase 1 & 4 & $-1/5$ & 0 & $2/5$ & 0 & 0 & -1 & $-6/5$ & 0 \\ 
Fase 2 & -10000 & -250 & 0 & -100 & 0 & -100 & 0 & 0 & 0 \\ 
\hline
$x_4$ & 52  & $22/5$ & 0 & $16/5$ & 1 & 1 & 0 & $-3/5$ & 0\\ 
$x_2$ & 16  & $1/5$ & 1 & $3/5$ & 0 & 0 & 0 & $1/5$ & 0\\ 
$a_3$ & 4  & $-1/5$ & 0 & $2/5$ & 0 & 0 & -1 & $-1/5$ & 1\\ 
\hline
\end{tabular}
\end{center}


\begin{enumerate}

\item Obtener la solución óptima correspondiente al plan de producción óptimo e interpretar la información que proporciona.

\item Dentro de qué rango de valores del margen de la aleación blanda no cambia la gama de producción en el plan óptimo de producción.

\item Dentro de qué rango de valores del precio de venta del mineral rojo, no cambia la gama de producción en el plan óptimo de producción.

\item Existe la posibilidad de producir una nueva aleación de alta gama, que requiere \textbf{6} toneladas de mineral rojo y \textbf{4} de mineral negro por cada tonelada de aleación. ¿Cuál es el beneficio unitario por tonelada de la nueva aleación que haría rentable su producción y venta? Esta nueva aleación no computaría para alcanzar la cantidad necesaria del acuerdo comercial.

\item Explicar el significado del $V^B_1$ en incluyendo explícitamente el significado de las tasas de sustitución de las variables $x_1$ con respecto a las variables básicas de la solución óptima.


\end{enumerate}


% PREGUNTAS ANTIGUAS

% \item \textbf Obtener la solución óptima del problema e {interpretar} el resultado y explicar el plan de
% extracción y el significado de los valores de las variables de holgura.


% \item Dado que el mineral negro se debe emplear completo, \textbf{justificar} si la empresa estaría interesada en extraer más o menos cantidad que las 80 toneladas semanales.

% \item \textbf{Justificar haciendo uso de los fundamentos del método del simplex} por qué, sin necesidad de resolver el problema, se podría afirmar que en la solución óptima $h_1$ nunca puede tomar un valor diferente de 0.

% \item Cuál es el intervalo relativo a la cantidad del compromiso comercial (aleaciones dura y fuerte) dentro del cual se mantiene la gama de producción correspondiente a la solución óptima.

% \item La empresa está interesada en saber cuál sería el impacto sobre su plan de producción y sobre su beneficio si, una determinada semana, no pudiera vender (y, por lo tanto, no produciría) más de 8 toneladas de aleación dura.
